\documentclass[11pt]{article}
\usepackage[T1]{fontenc}
\usepackage[utf8]{inputenc}
\usepackage{lmodern}
\usepackage[margin=1in]{geometry}
\usepackage{amsmath,amssymb,amsthm,mathtools}
\usepackage{microtype}
\usepackage{xcolor}
\usepackage{hyperref}
\hypersetup{colorlinks=true,linkcolor=blue!45!black,citecolor=blue!45!black,
  urlcolor=blue!45!black,pdftitle={Jacobians and Abel--Jacobi Maps}}
\setlength{\emergencystretch}{2em}
\numberwithin{equation}{section}
\newtheorem{theorem}{Theorem}[section]
\newtheorem{proposition}[theorem]{Proposition}
\newtheorem{lemma}[theorem]{Lemma}
\newtheorem{corollary}[theorem]{Corollary}
\theoremstyle{definition}
\newtheorem{definition}[theorem]{Definition}
\newtheorem{example}[theorem]{Example}
\newcommand{\Z}{\mathbf Z}
\newcommand{\R}{\mathbf R}
\newcommand{\C}{\mathbf C}
\newcommand{\Q}{\mathbf Q}
\newcommand{\AJ}{\operatorname{AJ}}
\newcommand{\End}{\operatorname{End}}
\newcommand{\Hom}{\operatorname{Hom}}
\newcommand{\Aut}{\operatorname{Aut}}
\newcommand{\ab}{\mathrm{ab}}
\newcommand{\dR}{\mathrm{dR}}
\newcommand{\HH}{\mathcal H}
\newcommand{\trans}{\mathsf T}

\title{Jacobians and Abel--Jacobi Maps}
\author{}
\date{}

\begin{document}
\maketitle
\vspace{-2em}

Let $X$ be a compact connected Riemann surface of genus $g\geq0$, with
base point $x_0$ and its complex orientation. Homology is singular
homology, $U^\vee=\Hom_{\C}(U,\C)$, and homomorphisms of complex tori
are holomorphic group homomorphisms. Put
\[
 W=H^0(X,\Omega_X^1),\qquad V=W^\vee.
\]

\section{Homology}

\begin{proposition}\label{prop:homology}
There is a presentation
\begin{equation}\label{eq:presentation}
 \pi_1(X,x_0)\cong
 \left\langle a_1,b_1,\ldots,a_g,b_g\ \middle|\
   \prod_{j=1}^g[a_j,b_j]=1\right\rangle,
 \qquad [a,b]=aba^{-1}b^{-1}.
\end{equation}
The Hurewicz map $h$ induces a canonical isomorphism
\[
 \pi_1(X,x_0)^{\ab}\xrightarrow{\sim}H_1(X,\Z),
 \qquad H_1(X,\Z)\cong\Z^{2g}.
\]
The homology classes of the generators form a symplectic basis:
\begin{equation}\label{eq:symplectic}
 a_i\cdot a_j=b_i\cdot b_j=0,\qquad a_i\cdot b_j=\delta_{ij}.
\end{equation}
Moreover, $\pi_1(X,x_0)$ is trivial for $g=0$, is $\Z^2$ for $g=1$,
and is nonabelian for $g\geq2$.
\end{proposition}

\begin{proof}
The polygon decomposition has one vertex, $2g$ edges, and one face
attached by $\prod_j[a_j,b_j]$; van Kampen gives
\eqref{eq:presentation}. For $g=0$ use the sphere with one vertex and
one face. If $F$ is the free group on the edges and $N$ is the normal
closure of the relator, then $N\subseteq[F,F]$, so
\[
 (F/N)^{\ab}=F/(N[F,F])=F/[F,F]\cong\Z^{2g}.
\]
Degree-one Hurewicz identifies this quotient with $H_1(X,\Z)$.
The standard oriented edge cycles have intersections
\eqref{eq:symplectic}. For $g=1$ the relation is commutativity; for
$g\geq2$ the quotient obtained by setting every $b_j=1$ is the
nonabelian free group on $a_1,\ldots,a_g$.
\end{proof}

\section{Periods}

Choose a compatible Riemannian metric and use
$*dx=dy$, $*dy=-dx$ in a holomorphic coordinate $z=x+iy$. Set
\[
 \HH^1(X,\R)=\{\alpha\in\Omega^1(X,\R):d\alpha=d*\alpha=0\}.
\]

\begin{theorem}[Hodge]
Every real de Rham class has a unique harmonic representative; thus
$\HH^1(X,\R)\xrightarrow{\sim}H^1_{\dR}(X,\R)$.
\end{theorem}

\begin{lemma}\label{lem:realparts}
The real-linear map
\[
 W_{\R}\xrightarrow{\sim}\HH^1(X,\R),\qquad
 \omega\longmapsto\operatorname{Re}\omega
\]
has inverse $\alpha\mapsto\alpha+i*\alpha$. Consequently,
\begin{equation}\label{eq:hodge}
 \dim_{\C}W=g,\qquad
 H^1_{\dR}(X,\C)\cong W\oplus\overline W.
\end{equation}
\end{lemma}

\begin{proof}
For $\alpha=p\,dx+q\,dy$,
\[
 \alpha+i*\alpha=(p-iq)\,dz,
 \qquad d\alpha=d*\alpha=0
 \iff q_x-p_y=p_x+q_y=0.
\]
These are the Cauchy--Riemann equations for $p-iq$. Conversely, the real
part of a holomorphic one-form satisfies these equations and recovers
the form by $\alpha+i*\alpha$. Hodge theory and de Rham's theorem give
$2\dim_{\C}W=\dim_{\R}H^1(X,\R)=2g$; complexifying the harmonic
representatives gives the decomposition in \eqref{eq:hodge}.
\end{proof}

Holomorphic one-forms are closed. Integration therefore defines
\begin{equation}\label{eq:period}
 P:H_1(X,\Z)\longrightarrow V,\qquad
 P(c)(\omega)=\int_c\omega,
\end{equation}
and its real-linear extension
$P_{\R}:H_1(X,\R)\to V_{\R}$.

\begin{definition}
A \emph{full lattice} in a real vector space is the integral span of a
real basis. A \emph{complex torus} is a quotient $U/\Gamma$, where $U$
is a finite-dimensional complex vector space and $\Gamma$ is a full
lattice in $U_{\R}$.
\end{definition}

\begin{theorem}\label{thm:lattice}
The map $P_{\R}$ is an isomorphism. Hence $P$ is injective and
\[
 \Lambda_X=P(H_1(X,\Z))
\]
is a full lattice of rank $2g$ in $V$.
\end{theorem}

\begin{proof}
If $P_{\R}(c)=0$, then
\[
 \int_c\operatorname{Re}\omega
 =\operatorname{Re}\int_c\omega=0\qquad(\omega\in W).
\]
By Lemma~\ref{lem:realparts} and Hodge theory, these real parts represent
all of $H^1_{\dR}(X,\R)$. Nondegeneracy of the de Rham pairing gives
$c=0$. Both spaces have real dimension $2g$, so $P_{\R}$ is an
isomorphism. If $c_1,\ldots,c_{2g}$ is an integral homology basis, then
$P(c_1),\ldots,P(c_{2g})$ is a real basis of $V_{\R}$ and
\[
 \Lambda_X=\bigoplus_{j=1}^{2g}\Z P(c_j).
\]
Thus $\Lambda_X$ is discrete, and the compact parallelepiped
$\{\sum_jt_jP(c_j):0\leq t_j\leq1\}$ surjects onto $V/\Lambda_X$.
\end{proof}

\begin{definition}
The \emph{Jacobian} of $X$ is
\begin{equation}\label{eq:jacobian}
 J(X)=V/\Lambda_X
 =H^0(X,\Omega_X^1)^\vee/P(H_1(X,\Z)).
\end{equation}
\end{definition}

Lattice translations act freely and properly discontinuously by
biholomorphisms. Hence $J(X)$ is a compact complex Lie group of complex
dimension $g$, and $J(X)\cong\R^{2g}/\Z^{2g}$ as real manifolds after
choosing an integral homology basis.

\section{Abel--Jacobi Maps}

\begin{definition}
For $x\in X$, choose a piecewise smooth path $c$ from $x_0$ to $x$ and set
\begin{equation}\label{eq:aj}
 \AJ_{x_0}(x)=\left[\omega\longmapsto\int_c\omega\right]\in J(X).
\end{equation}
\end{definition}

\begin{proposition}\label{prop:aj}
The map $\AJ_{x_0}$ is well defined and holomorphic, with
$\AJ_{x_0}(x_0)=0$. Under $T_{\AJ_{x_0}(x)}J(X)=V$,
\begin{equation}\label{eq:daj}
 d(\AJ_{x_0})_x(v)(\omega)=\omega_x(v).
\end{equation}
For another base point $y_0$,
\[
 \AJ_{y_0}(x)=\AJ_{x_0}(x)-\AJ_{x_0}(y_0).
\]
\end{proposition}

\begin{proof}
Two choices of $c$ differ by an integral cycle, so their integral
functionals differ by an element of $\Lambda_X$. On a coordinate disk,
holomorphic primitives give a holomorphic local lift to $V$ and imply
\eqref{eq:daj}. Concatenation of paths gives the base-point formula.
\end{proof}

\begin{lemma}\label{lem:cover}
For a complex torus $T=U/\Gamma$, the quotient $q:U\to T$ is a universal
cover. Lifting based loops from $0$ defines an isomorphism
\[
 \delta_T:\pi_1(T,0)\xrightarrow{\sim}\Gamma,\qquad
 [\gamma]\longmapsto\widetilde\gamma(1).
\]
\end{lemma}

\begin{proof}
The lattice action gives a covering and $U$ is contractible. Endpoints
of lifted concatenations add. Each $\lambda\in\Gamma$ is realized by
$t\mapsto q(t\lambda)$, and a lift with endpoint $0$ is a loop in the
simply connected space $U$.
\end{proof}

\begin{theorem}\label{thm:abelianization}
Writing $\delta=\delta_{J(X)}$, one has
\begin{equation}\label{eq:main}
 \delta\circ(\AJ_{x_0})_*=P\circ h.
\end{equation}
Thus $(\AJ_{x_0})_*$ is surjective and
\[
 \ker((\AJ_{x_0})_*)=[\pi_1(X,x_0),\pi_1(X,x_0)],\qquad
 \pi_1(J(X),0)\cong\pi_1(X,x_0)^{\ab}.
\]
Moreover, $\AJ_{x_0}$ induces an isomorphism
$H_1(X,\Z)\xrightarrow{\sim}H_1(J(X),\Z)$.
\end{theorem}

\begin{proof}
For a based loop $\gamma:[0,1]\to X$, the lift of
$\AJ_{x_0}\circ\gamma$ starting at zero is
\[
 \widetilde\gamma(t)(\omega)=\int_{\gamma|_{[0,t]}}\omega.
\]
Its endpoint is $P(h([\gamma]))$, proving \eqref{eq:main}.
The assertions about fundamental groups follow from injectivity of $P$
and degree-one Hurewicz. Abelianizing \eqref{eq:main} and using naturality
of Hurewicz proves the homology assertion.
\end{proof}

\section{Universal Property}

\begin{lemma}\label{lem:linear}
For complex tori $T_i=U_i/\Gamma_i$,
\[
 \Hom(T_1,T_2)\cong
 \{A\in\Hom_{\C}(U_1,U_2):A(\Gamma_1)\subseteq\Gamma_2\}.
\]
Also, $H^0(T_i,\Omega_{T_i}^1)=U_i^\vee$ via invariant one-forms.
\end{lemma}

\begin{proof}
A homomorphism lifts uniquely to a holomorphic $A:U_1\to U_2$ with
$A(0)=0$. The continuous map
$(u,v)\mapsto A(u+v)-A(u)-A(v)$ takes values in $\Gamma_2$ and vanishes
at $(0,0)$, so $A$ is additive. Hence $dA_u=dA_0$ and $A(u)=dA_0(u)$
is complex-linear. It maps $\Gamma_1$ into $\Gamma_2$. Conversely, such
a linear map descends to a homomorphism. The coefficients of the lift
of a holomorphic one-form to $U_i$ are entire and $\Gamma_i$-periodic;
they are bounded and thus constant by Liouville's theorem.
\end{proof}

\begin{theorem}\label{thm:universal}
Under $H^0(J(X),\Omega_{J(X)}^1)=V^\vee=W$, the map
\[
 \AJ_{x_0}^*:H^0(J(X),\Omega_{J(X)}^1)\xrightarrow{\sim}W
\]
is the identity. For every complex torus $T$ and holomorphic map
$f:X\to T$, there is a unique homomorphism $F:J(X)\to T$ such that
\begin{equation}\label{eq:universal}
 f(x)=f(x_0)+F(\AJ_{x_0}(x)).
\end{equation}
\end{theorem}

\begin{proof}
The assertion on one-forms follows from \eqref{eq:daj}.
After translating $f$, assume $f(x_0)=0$ and write $T=U/\Gamma$.
Dualizing pullback of invariant forms gives
\[
 A=(f^*)^\vee:V\longrightarrow U,\qquad
 \ell(A(v))=v(f^*\ell)\quad(\ell\in U^\vee).
\]
For a path $c$ from $x_0$, let $u_c$ be the endpoint of the lift of
$f\circ c$ to $U$ starting at zero. Then
\[
 \ell\left(A\left(\omega\longmapsto\int_c\omega\right)\right)
 =\int_c f^*\ell=\ell(u_c).
\]
For a loop, $u_c\in\Gamma$. Since $h$ is surjective, this implies
$A(\Lambda_X)\subseteq\Gamma$. Lemma~\ref{lem:linear} gives a
homomorphism $F:J(X)\to T$, and the same identity for an arbitrary
path proves \eqref{eq:universal}.
If $F_1,F_2$ give the factorization, their induced maps on fundamental
groups agree after composition with $(\AJ_{x_0})_*$. This map is
surjective by Theorem~\ref{thm:abelianization}, so the linear lifts of
$F_1,F_2$ agree on $\Lambda_X$. The lattice spans $V_{\R}$, hence the
lifts and the homomorphisms agree.
\end{proof}

\section{Polarization}

Assume $g\geq1$ and choose the symplectic basis
\eqref{eq:symplectic}.

\begin{lemma}\label{lem:bilinear}
For closed complex one-forms $\alpha,\beta$,
\begin{equation}\label{eq:bilinear}
 \int_X\alpha\wedge\beta
 =\sum_{j=1}^g\left(
 \int_{a_j}\alpha\int_{b_j}\beta
 -\int_{b_j}\alpha\int_{a_j}\beta\right).
\end{equation}
For $0\neq\omega\in W$,
\begin{equation}\label{eq:positivity}
 \frac{i}{2}\int_X\omega\wedge\overline\omega>0.
\end{equation}
\end{lemma}

\begin{proof}
If $u_j,v_j$ are the cohomology classes dual to $a_j,b_j$, Poincar\'e
duality gives $\int_Xu_i\smile u_j=\int_Xv_i\smile v_j=0$ and
$\int_Xu_i\smile v_j=\delta_{ij}$. Expanding $[\alpha],[\beta]$ gives
\eqref{eq:bilinear}. Locally,
$\frac{i}{2}\omega\wedge\overline\omega=|f(z)|^2dx\wedge dy$
for $\omega=f(z)\,dz$, proving \eqref{eq:positivity}.
\end{proof}

\begin{proposition}\label{prop:periodmatrix}
There is a unique basis $\omega_1,\ldots,\omega_g$ of $W$ satisfying
$\int_{a_j}\omega_i=\delta_{ij}$. For
\[
 \tau_{ij}=\int_{b_j}\omega_i,\qquad Y=\operatorname{Im}\tau,
\]
one has $\tau^{\trans}=\tau$ and $Y>0$. In the coordinates
$v\mapsto(v(\omega_1),\ldots,v(\omega_g))^{\trans}$,
\begin{equation}\label{eq:normalized}
 \Lambda_X=\Z^g+\tau\Z^g,\qquad
 J(X)\cong\C^g/(\Z^g+\tau\Z^g).
\end{equation}
\end{proposition}

\begin{proof}
If all $a$-periods of $\omega$ vanish, \eqref{eq:bilinear} gives
$\int_X\omega\wedge\overline\omega=0$; thus $\omega=0$ by
\eqref{eq:positivity}. Since $\dim_{\C}W=g$, the $a$-period map
$W\to\C^g$ is an isomorphism. This gives the normalized basis and
\eqref{eq:normalized}. Since $\omega_i\wedge\omega_j=0$,
\eqref{eq:bilinear} gives $0=\tau_{ji}-\tau_{ij}$. For
$\omega=\sum_i c_i\omega_i$ it gives
\[
 \frac{i}{2}\int_X\omega\wedge\overline\omega
 =\overline c^{\trans}Yc>0\qquad(c\neq0),
\]
so $Y$ is positive definite.
\end{proof}

\begin{definition}
A \emph{Riemann form} on a complex torus $U/\Gamma$ is an integral
alternating form $E$ on $\Gamma$ whose real-bilinear extension satisfies
\[
 E(iu,iv)=E(u,v),\qquad E(u,iu)>0\quad(u\neq0).
\]
It defines a polarization, which is \emph{principal} if
$\Gamma\to\Hom_{\Z}(\Gamma,\Z)$, $\lambda\mapsto E(\lambda,-)$,
is an isomorphism.
\end{definition}

\begin{theorem}[Riemann's criterion]
A complex torus is projective if and only if it admits a Riemann form.
A projective complex torus is an abelian variety over $\C$.
\end{theorem}

\begin{theorem}\label{thm:polarization}
The form
\[
 E(z,w)=\operatorname{Im}\bigl(\overline z^{\trans}Y^{-1}w\bigr)
\]
defines a canonical principal polarization on $J(X)$, characterized by
\begin{equation}\label{eq:intersection}
 E(P(c),P(d))=c\cdot d\qquad(c,d\in H_1(X,\Z)).
\end{equation}
Thus $J(X)$ is a principally polarized abelian variety.
\end{theorem}

\begin{proof}
The positive definite Hermitian form
$H(z,w)=\overline z^{\trans}Y^{-1}w$ is conjugate-linear in its first
variable. Hence $E$ is alternating, $E(iz,iw)=E(z,w)$, and
$E(z,iz)=H(z,z)>0$ for $z\neq0$. Writing $\tau=B+iY$, with
$B^{\trans}=B$, gives
\begin{align*}
 E(m+\tau n,m'+\tau n')
 &=(m+Bn)^{\trans}n'-n^{\trans}(m'+Bn')\\
 &=m^{\trans}n'-n^{\trans}m'
 \qquad(m,n,m',n'\in\Z^g).
\end{align*}
Thus $E$ is integral and has matrix
$\mathsf S=\left(\begin{smallmatrix}0&I_g\\-I_g&0\end{smallmatrix}\right)$
in the period basis. This matrix is unimodular, so the polarization is
principal. The same computation gives \eqref{eq:intersection}, which
also proves independence of the symplectic basis. Riemann's criterion
gives projectivity.
\end{proof}

\section{Endomorphisms}

Let $T=U/\Gamma$ have complex dimension $g$, and let $\mathcal I$ denote
multiplication by $i$ on $U_{\R}=\Gamma\otimes_{\Z}\R$.

\begin{proposition}\label{prop:endomorphisms}
Restriction to $\Gamma$ identifies the endomorphism ring with
\begin{align*}
 \End(T)
 &\cong\{A\in\End_{\C}(U):A(\Gamma)\subseteq\Gamma\}\\
 &\cong\{u\in\End_{\Z}(\Gamma):
       u_{\R}\mathcal I=\mathcal Iu_{\R}\}
 \ \subseteq\ \End_{\Z}(\Gamma).
\end{align*}
After choosing a lattice basis, the last ring is $M_{2g}(\Z)$.
In particular, $\End(T)$ is a free abelian group of finite rank at most
$4g^2$.
\end{proposition}

\begin{proof}
The first description is Lemma~\ref{lem:linear}. A real-linear map is
complex-linear exactly when it commutes with $\mathcal I$. Since
$\Gamma$ spans $U_{\R}$, restriction is injective. A subgroup of the
free abelian group $M_{2g}(\Z)$ is free of rank at most $4g^2$.
\end{proof}

\begin{proposition}
For a principally polarized torus $(T,E)$, its polarization-preserving
automorphisms correspond, in a symplectic lattice basis, to
\[
 \Aut(T,E)\cong
 \{M\in M_{2g}(\Z):M_{\R}\mathcal I=\mathcal I M_{\R},\
                  M^{\trans}\mathsf S M=\mathsf S\}.
\]
For each $n\in\Z$, multiplication by $n$ belongs to $\End(T)$ and
satisfies $E(n\lambda,n\mu)=n^2E(\lambda,\mu)$.
\end{proposition}

\begin{proof}
The first equation is the complex-linearity condition in
Proposition~\ref{prop:endomorphisms}; the second is preservation of $E$.
The latter gives an integral inverse
$M^{-1}=\mathsf S^{-1}M^{\trans}\mathsf S$, so the induced map is an
automorphism. The assertion about multiplication follows from
bilinearity.
\end{proof}

\begin{example}\label{ex:elliptic}
For $\operatorname{Im}\tau>0$, set $E_\tau=\C/(\Z+\tau\Z)$. Then
\[
 \End(E_\tau)=\{\alpha\in\C:
       \alpha(\Z+\tau\Z)\subseteq\Z+\tau\Z\},
\]
and
\[
 \End(E_\tau)\neq\Z
 \iff \tau\text{ is imaginary quadratic over }\Q.
\]
In particular,
\[
 \End(E_i)=\Z[i],\qquad \End(E_{\,i\sqrt[4]{2}})=\Z,
 \qquad \pi_1(E_i)\cong\pi_1(E_{\,i\sqrt[4]{2}})\cong\Z^2.
\]
\end{example}

\begin{proof}
A complex-linear endomorphism of $\C$ is multiplication by a scalar
$\alpha$. The lattice condition is equivalent to
\[
 \alpha=a+b\tau,\qquad \alpha\tau=c+d\tau,
 \qquad a,b,c,d\in\Z,
\]
and hence to
\begin{equation}\label{eq:quadratic}
 b\tau^2+(a-d)\tau-c=0.
\end{equation}
If $\tau$ is not imaginary quadratic, the polynomial must vanish
identically: $b=c=0$, $a=d$, so $\alpha\in\Z$. Conversely, if
$A\tau^2+B\tau+C=0$ with integers $A\neq0,B,C$, then multiplication
by $A\tau$ preserves the lattice and is not integer multiplication.
Thus the exceptional parameters form a countable set.
For $\tau=i$, the lattice is $\Z[i]$ and its scalar multiplier ring is
$\Z[i]$. For $\tau=i\sqrt[4]{2}$, \eqref{eq:quadratic} gives
$a=d$ and $b\sqrt2+c=0$, so $b=c=0$.
\end{proof}

\section{Low Genus}

\begin{proposition}
If $g=0$, then $J(X)$ is a point. If $g=1$, then
$\AJ_{x_0}:X\to J(X)$ is a biholomorphism. If $g\geq2$, then
$\AJ_{x_0}$ is not a homotopy equivalence.
\end{proposition}

\begin{proof}
For $g=0$, $W=0$. For $g=1$, Theorem~\ref{thm:universal} implies that
$\AJ_{x_0}$ is nonconstant, so it is a surjective holomorphic map of
positive degree. By Theorem~\ref{thm:abelianization}, its pullback on
integral $H^1$ is an isomorphism. Choose $\alpha,\beta\in H^1(J(X),\Z)$
with $\langle\alpha\smile\beta,[J(X)]\rangle=1$. Their pullbacks form
an integral basis of $H^1(X,\Z)$, so
\[
 \deg(\AJ_{x_0})
 =\left\langle\AJ_{x_0}^*\alpha\smile\AJ_{x_0}^*\beta,[X]\right\rangle
 =\pm1.
\]
Positivity gives degree one. Positive local multiplicities then imply
that the map is bijective and unramified, hence biholomorphic.
For $g\geq2$, $\pi_1(X,x_0)$ is nonabelian whereas
$\pi_1(J(X),0)\cong\Z^{2g}$.
\end{proof}

For $X=E_\tau$ and $x_0=[0]$, the form $dz$ identifies $V$ with $\C$;
its period lattice is $\Z+\tau\Z$ and $\AJ_0([z])=[z]$.

\begin{thebibliography}{9}
\bibitem{hatcher}
A. Hatcher, \emph{Algebraic Topology}, Cambridge University Press, 2002.
\href{https://pi.math.cornell.edu/~hatcher/AT/ATpage.html}{Online edition}.

\bibitem{forster}
O. Forster, \emph{Lectures on Riemann Surfaces}, translated by B. Gilligan,
Graduate Texts in Mathematics 81, Springer, 1981.

\bibitem{milneav}
J. S. Milne, \emph{Abelian Varieties}, course notes, version 2.00, 2008.
\href{https://www.jmilne.org/math/CourseNotes/AV.pdf}{Online notes}.

\bibitem{milnejv}
J. S. Milne, \emph{Jacobian Varieties}, in \emph{Arithmetic Geometry},
G. Cornell and J. H. Silverman, eds., Springer, 1986, pp.~167--212.
\href{https://www.jmilne.org/math/xnotes/JVs.html}{Revised online edition}.
\end{thebibliography}

\end{document}
